\SourceLeaf{046}
Ita e. g. pro $m=2$ habemus
\[
p=-{l'}^2+4l''
\]
Perinde pro $m=3$ invenitur
\[
p=-{l'}^2{l''}^2+4{l'}^3l'''+4{l''}^3-18l'l''l'''+27{l'''}^2
\]
Determinans functionis $y$ itaque est functio coëfficientium $l',l'',l'''$ etc. talis,
quae per substitutiones $l'=\lambda',l''=\lambda'',l'''=\lambda'''$ etc. transit in productum ex
omnibus differentiis inter binas quantitatum $a,b,c$ etc. In casu eo, ubi $m=1$,
i. e. ubi unica tantum indeterminata $a$ habetur, adeoque nullae omnino adsunt
differentiae, ipsum numerum 1 tamquam determinantem functionis $y$ adoptare
conveniet.

In stabilienda notione determinantis, coëfficientes functionis $y$ tamquam
quantitates indeterminatas spectare oportuit. Determinans functionis cum coëfficientibus determinatis
\[
Y=x^m-L'x^{m-1}+L''x^{m-2}-L'''x^{m-3}+\text{etc.}
\]
erit numerus determinatus $P$, puta valor functionis $p$ pro $l'=L',l''=L''$,
$l'''=L'''$ etc. Quodsi itaque supponimus, $Y$ resolvi posse in factores simplices
\[
Y=(x-A)(x-B)(x-C)\, .\, .\, .\, .\, .
\]
sive $Y$ oriri ex
\[
\upsilon=(x-a)(x-b)(x-c)\, .\, .\, .\, .\, .
\]
statuendo $a=A,b=B,c=C$ etc., adeoque per easdem substitutiones $\lambda',\lambda'',\lambda'''$ etc.
resp. fieri $L',L'',L'''$ etc., manifesto $P$ aequalis erit producto e factoribus
\[
\begin{aligned}
&(A-B)(A-C)(A-D)\, .\, .\, .\, .\, .\\
\times{}&(B-A)(B-C)(B-D)\, .\, .\, .\, .\, .\\
\times{}&(C-A)(C-B)(C-D)\, .\, .\, .\, .\, .\\
\times{}&(D-A)(D-B)(D-C)\, .\, .\, .\, .\, .\\
&\text{etc.}
\end{aligned}
\]
Patet itaque, si fiat $P=0$, inter quantitates $A,B,C$ etc. duas saltem aequales reperiri debere; contra, si non fuerit $P=0$, cunctas $A,B,C$ etc. necessario inaequales esse. Iam observamus, si statuamus $\frac{\mathrm dY}{\mathrm dx}=Y'$, sive

\SourceLeaf{047}
\[
\begin{gathered}Y'=mx^{m-1}-(m-1)L'x^{m-2}+(m-2)L''x^{m-3}\\-(m-3)L'''x^{m-4}+\text{etc.}\end{gathered}
\]
haberi
\[
\begin{aligned}
Y'={}&(x-B)(x-C)(x-D)\, .\, .\, .\, .\, .\\
&+(x-A)(x-C)(x-D)\, .\, .\, .\, .\, .\\
&+(x-A)(x-B)(x-D)\, .\, .\, .\, .\, .\\
&+(x-A)(x-B)(x-C)\, .\, .\, .\, .\, .\\
&+\text{etc.}
\end{aligned}
\]
Si itaque duae quantitatum $A,B,C$ etc. aequales sunt, e. g. $A=B$, $Y'$ per
$x-A$ divisibilis erit, sive $Y$ et $Y'$ implicabunt divisorem communem $x-A$.
Vice versa, si $Y'$ cum $Y$ ullum divisorem communem habere supponitur, necessario $Y'$ aliquem factorem simplicem ex his $x-A,x-B,x-C$ etc. implicare
debebit, e. g. primum $x-A$, quod manifesto fieri nequit, nisi $A$ alicui reliquarum $B,C,D$ etc. aequalis fuerit. Ex his omnibus itaque colligimus duo
\textsc{Theoremata}:

I. \textit{Si determinans functionis $Y$ fit $=0$, certo $Y$ cum $Y'$ divisorem communem habet, adeoque, si $Y$ et $Y'$ divisorem communem non habent, determinans functionis $Y$ nequit esse $=0$.}

II. \textit{Si determinans functionis $Y$ non est $=0$, certo $Y$ et $Y'$ divisorem communem habere nequeunt; vel, si $Y$ et $Y'$ divisorem communem habent, necessario determinans functionis $Y$ esse debet $=0$.}

\SecondArticle{7}
At probe notandum est, totam vim huius demonstrationis simplicissimae inniti suppositioni, functionem $Y$ in factores simplices resolvi posse: quae ipsa suppositio, hocce quidem loco, ubi de demonstratione generali huius resolubilitatis
agitur, nihil esset nisi petitio principii. Et tamen a paralogismis huic prorsus similibus non sibi caverunt omnes, qui demonstrationes analyticas theorematis principalis tentaverunt, cuius speciosae illusionis originem iam in ipsa disquisitionis
enunciatione animadvertimus, quum omnes in \textit{formam} tantum radicum aequationum inquisiverint, dum \textit{existentiam} temere suppositam demonstrare oportuisset.
Sed de tali procedendi modo, qui nimis a rigore et claritate abhorret, satis iam in
commentatione supra citata dictum est. Quamobrem iam theoremata art. praec.,
\SourceLeaf{048}
quorum altero saltem ad propositum nostrum non possumus carere, solidiori fundamento superstruemus: a secundo, tamquam faciliori initium faciemus.

\SecondArticle{8}
Denotemus per $\rho$ functionem
\[
\begin{aligned}
&\frac{\pi(x-b)(x-c)(x-d)\, .\, .\, .\, .\, .}{(a-b)^2(a-c)^2(a-d)^2\, .\, .\, .\, .\, .}\\
&+\frac{\pi(x-a)(x-c)(x-d)\, .\, .\, .\, .\, .}{(b-a)^2(b-c)^2(b-d)^2\, .\, .\, .\, .\, .}\\
&+\frac{\pi(x-a)(x-b)(x-d)\, .\, .\, .\, .\, .}{(c-a)^2(c-b)^2(c-d)^2\, .\, .\, .\, .\, .}\\
&+\frac{\pi(x-a)(x-b)(x-c)\, .\, .\, .\, .\, .}{(d-a)^2(d-b)^2(d-c)^2\, .\, .\, .\, .\, .}\\
&+\text{etc.}
\end{aligned}
\]
quae, quoniam $\pi$ per singulos denominatores est divisibilis, fit functio integra
indeterminatarum $x,a,b,c$ etc. Statuamus porro $\frac{\mathrm d\upsilon}{\mathrm dx}=\upsilon'$, ita ut habeatur
\[
\begin{aligned}
\upsilon'={}&(x-b)(x-c)(x-d)\, .\, .\, .\, .\, .\\
&+(x-a)(x-c)(x-d)\, .\, .\, .\, .\, .\\
&+(x-a)(x-b)(x-d)\, .\, .\, .\, .\, .\\
&+(x-a)(x-b)(x-c)\, .\, .\, .\, .\, .\\
&+\text{etc.}
\end{aligned}
\]
Manifesto pro $x=a$, fit $\rho\upsilon'=\pi$, unde concludimus, functionem $\pi-\rho\upsilon'$ indefinite divisibilem esse per $x-a$, et perinde per $x-b,x-c$ etc., nec non
per productum $\upsilon$. Statuendo itaque
\[
\frac{\pi-\rho\upsilon'}{\upsilon}=\sigma
\]
erit $\sigma$ functio integra indeterminatarum $x,a,b,c$ etc., et quidem, perinde ut $\rho$,
symmetrica ratione indeterminatarum $a,b,c$ etc. Erui poterunt itaque functiones duae integrae $r,s$, indeterminatarum $x,l',l'',l'''$ etc., tales quae per substitutiones $l'=\lambda',l''=\lambda'',l'''=\lambda'''$ etc. transeant in $\rho,\sigma$ resp. Quodsi itaque analogiam sequentes, functionem
\[
\begin{gathered}mx^{m-1}-(m-1)l'x^{m-2}+(m-2)l''x^{m-3}\\-(m-3)l'''x^{m-4}+\text{etc.}\end{gathered}
\]
i. e. quotientem differentialem $\frac{\mathrm dy}{\mathrm dx}$ per $y'$ denotemus, ita ut $y'$ per easdem illas
\SourceLeaf{049}
substitutiones transeat in $\upsilon'$. patet, $p-sy-ry'$ per easdem substitutiones transire
in $\pi-\sigma\upsilon-\rho\upsilon'$, i. e. in 0, adeoque necessario iam per se identice evanescere
debere (art. 5): habemus proin aequationem identicam
\[
p=sy+ry'
\]
Hinc si supponamus, ex substitutione $l'=L',l''=L'',l'''=L'''$ etc. prodire
$r=R,s=S$, erit etiam identice
\[
P=SY+RY'
\]
ubi quum $S,R$ sint functionis integrae ipsius $x$, $P$ vero quantitas determinata
seu numerus, sponte patet, $Y$ et $Y'$ divisorem communem habere non posse, nisi
fuerit $P=0$. Quod est ipsum theorema posterius art. 6.

\SecondArticle{9}
Demonstrationem theorematis prioris ita absolvemus, ut ostendamus, in
casu eo, ubi $Y$ et $Y'$ non habent divisorem communem, certo fieri non posse
$P=0$. Ad hunc finem primo, per praecepta art. 2 erutas supponimus duas
functiones integras indeterminatae $x$, puta $fx$ et $\varphi x$, tales, ut habeatur aequatio identica
\[
fx\, .\,Y+\varphi x\, .\,Y'=1
\]
quam hic ita exhibemus:
\[
fx\, .\,\upsilon+\varphi x\, .\,\upsilon'=1+fx\, .\,(\upsilon-Y)+\varphi x\, .\,\frac{\mathrm d(\upsilon-Y)}{\mathrm dx}
\]
sive, quoniam habemus
\[
\begin{aligned}
\upsilon'={}&(x-b)(x-c)(x-d)\, .\, .\, .\, .\\
&+(x-a)\, .\,\frac{\mathrm d[(x-b)(x-c)(x-d)\, .\, .\, .\, .]}{\mathrm dx}
\end{aligned}
\]
in forma sequente:
\[
\begin{aligned}
&\varphi x\, .\,(x-b)(x-c)(x-d)\, .\, .\, .\, .\\
&+\varphi x\, .\,(x-a)\, .\,\frac{\mathrm d[(x-b)(x-c)(x-d)\, .\, .\, .\, .]}{\mathrm dx}\\
&+fx\, .\,(x-a)(x-b)(x-c)(x-d)\, .\, .\, .\, .\\&=1+fx\, .\,(\upsilon-Y)+\varphi x\, .\,\frac{\mathrm d(\upsilon-Y)}{\mathrm dx}
\end{aligned}
\]
Exprimamus brevitatis caussa
\[
fx\, .\,(y-Y)+\varphi x\, .\,\frac{\mathrm d(y-Y)}{\mathrm dx}
\]
